\documentclass[withoutpreface,bwprint]{cumcmthesis} %去掉封面与编号页
\usepackage{url} % 支持网址链接，便于引用网络资源
\usepackage{cite} % 优化参考文献引用格式
\usepackage{graphicx} % 插入和管理图片
\usepackage{pdflscape} % 使页面横向显示，适合宽表格
\usepackage{lscape} % 旋转页面内容，适合横向表格或图片
\usepackage{tabularx} % 支持自动调整宽度的表格
\usepackage{afterpage} % 控制内容在下一页显示
\usepackage{longtable} % 支持跨页长表格

\usepackage[numbers,sort&compress]{natbib} % 参考文献排序与压缩显示

\usepackage[figuresright]{rotating} % 支持图表旋转


\usepackage{amsmath} % 数学公式增强
\usepackage{amssymb} % 更多数学符号
\usepackage{geometry} % 页面边距自定义
\usepackage{enumitem} % 列表样式自定义
\usepackage{caption} % 图表标题样式自定义
\usepackage{fancyhdr} % 页眉页脚自定义
\usepackage{titlesec} % 章节标题样式自定义

\usepackage{ctex} % 中文支持，适合中文文档

\usepackage{listings} % 插入和高亮显示代码
% 设置代码样式
\lstset{
	numbers             =   left,               % 行号显示在左侧
	showspaces          =   false,              % 不显示空格符号
	numberstyle         =   \zihao{-5}\ttfamily,% 行号风格，小五号，等宽字体
	showstringspaces    =   false,              % 不显示字符串中的空格
	captionpos          =   t,                  % 标题位置在顶部
	frame               =   lrtb,               % 显示四周边框
	breaklines          =   true,               % 自动换行
	columns             =   fixed,              % 固定列宽
	basewidth           =   0.5em,              % 基础宽度
	caption             =   \raggedright,       % 标题左对齐
}

% Python 代码样式
\lstdefinestyle{Python}{
	language        =   Python, % 语言设为Python
	basicstyle      =   \zihao{-5}\ttfamily\color{black}, % 基本风格，小五号，等宽字体，黑色
	keywordstyle    =   \color{blue}\bfseries\itshape, % 关键字风格，蓝色加粗斜体
	stringstyle     =   \color{magenta}, % 字符串风格，洋红色
	commentstyle    =   \color{red}\itshape, % 注释风格，红色斜体
}

% MATLAB 代码样式
\lstdefinestyle{MATLAB}{
	language        =   Matlab, % 语言设为MATLAB
	basicstyle      =   \zihao{-5}\ttfamily\color{black}, % 基本风格，小五号，等宽字体，黑色
	keywordstyle    =   \color{green}\bfseries\itshape, % 关键字风格，绿色加粗斜体
	stringstyle     =   \color{purple}, % 字符串风格，紫色
	commentstyle    =   \color{gray}\itshape, % 注释风格，绿色斜体
}



\usepackage{xcolor} % 颜色支持，代码和表格可用
\usepackage{multicol} % 多栏排版

\usepackage{booktabs} % 优美的表格线
\usepackage{hyperref} % 支持PDF跳转、超链接
\usepackage{subcaption} % 子图标题，多个图片并排显示
\usepackage {float}  % 更灵活地控制浮动体（如表格、图片）位置
\usepackage{array} % 增强表格功能
\hypersetup{hidelinks,
	colorlinks=true,
	allcolors=black,
	pdfstartview=Fit,
	breaklinks=true
}
\usepackage{siunitx} % 科学计数法和单位格式化

\title{基于主要模型对关键词的研究}
%\tihao{A}					%题号
%\baominghao{4321}			%报名号
%\schoolname{BIT}			%学校
%\yearinput{2025}			%年份
%\monthinput{09}			%月份
%\dayinput{03}				%日期

\begin{document}


\input{texfile/1abstract}%摘要  abstract.tex

\input{texfile/2ProblemRestatement}%插入问题重述   ProblemRestatement.tex

\input{texfile/3ProblemAnalysis}%插入问题分析    ProblemAnalysis.tex

\input{texfile/4AssumptionAndSign}%插入模型假设及符号说明  AssumptionAndSign.tex

\input{texfile/5MakeModel} %插入模型的建立与求解   MakeModel.tex

\input{texfile/6ErrorAnalysis} %插入误差分析   ErrorAnalysis.tex

\input{texfile/7ModelEvaluation} %插入模型评价	ModelEvaluation.tex

\newpage
%参考文献
%引用     \textsuperscript{\cite{ref01}}
\begin{thebibliography}{9}%宽度9
	
	\setlength{\itemsep}{-1mm}



\bibitem{ref01}


%\textsuperscript{\cite{ref01}}
\bibitem{ref02}


%\textsuperscript{\cite{ref02}}
\bibitem{ref03}


%\textsuperscript{\cite{ref03}}
\bibitem{ref04}


%\textsuperscript{\cite{ref04}}
\bibitem{ref05}


%\textsuperscript{\cite{ref05}}
\bibitem{ref06}


%\textsuperscript{\cite{ref06}}
\bibitem{ref07}

%\textsuperscript{\cite{ref07}}
\bibitem{ref08}


%\textsuperscript{\cite{ref08}}
\bibitem{ref09}


%\textsuperscript{\cite{ref09}}
\bibitem{ref10}


%\textsuperscript{\cite{ref10}}
	
	

	
\end{thebibliography}	 %插入参考文献   Reference.tex

\input{texfile/9Appendix} %插入附录   Appendix.tex



\end{document}